\section{Approach and Methodology}

For the approach and methodology, a sequential flow will be used on three models. The three models are 
SmolVLM-500M, BLIP and MiniCPIM-V 4.6. Each model has different capabilities and architectures.
First, the SFT baseline will be established. For each model, a conventional SFT condition will be trained
using full image + question and ground-truth answer. Thereafter, the OPSD fine-tuning will be performed. 
Two policies will be formed, one for the student and one for the teacher. The student receives the full image
and question and generates its own answer trajectory. The privileged teacher receives a relevant region (cropped portion of image) 
that contains more focused visual evidence relative to the question \parencite{yuan2026vision}.
The teacher then evaluates the student's generated trajectory and provides token-level probability distributions
\parencite{zhao2601self}. The student is trained to minimise the divergence between these two policies. 

This method could outperform SFT for two main arguments. Firstly, because of on-policy supervision, the train-inference
distribution mismatch will be reduced. Secondly, OPSD provides the student with the teacher's full 
next-token probability distribution which may be more useful than simply receiving one correct answer token.

